\section{Predictor family and feature realizations}
\label{sec:method}

\subsection{A compact offline scorer}
We fit a ridge-regularized logistic model by Newton--IRLS,
$f(\mathbf x)=\sigma(\mathbf w^\top\mathbf x+b)$; on the within-layer and
cross-layer features it has two coefficients and one intercept. A
log-loss fit without class reweighting has low measured ECE here and the
class-balanced LightGBM does not, but an unweighted or recalibrated tree removes
that difference (\S\ref{sec:exp}): calibration is separate from ranking here.

\subsection{Offline predictor versus runtime reconstruction}
\label{sec:realization}
The masking-loop scorer reconstructs the cross-layer view from the previous
layer's EMA, neutralizes the query view, and feeds selection into later states, so
its task F1 cannot inherit the offline predictor's AUC. Table~\ref{tab:protocol}
distinguishes fixed offline features, policy-dependent runtime features and the
irreversible physical reference.
\textbf{Bridge experiment.} With the frozen archived two-view checkpoint, one label
(top-$\lceil0.1n\rceil$ of block attention four steps ahead among the blocks present
now and still present then) and the 128 frozen prompts of \S\ref{sec:deploy}, we
score (A) offline-style features re-derived from the full-cache attention stream under
version-2 conventions, (B) the reference backend's runtime realization on the same
rows, and (C) that realization along the on-policy trajectory at 20\% budget. In A and B the candidates are every block present; in C they are the blocks
retained now \emph{and} at $t+h$, a survivor-conditioned arm whose $n$ and label
denominator are taken among those survivors. Mean per-request AUC is
0.902/0.899/0.989 on Llama-3.1-8B (bfloat16) and 0.894/0.892/0.989 on Qwen2.5-7B
(float32): the feature realization costs $-0.003$ $[-0.004,-0.001]$ (84/128
requests negative) and $-0.002$ $[-0.003,-0.000]$ (75/128), whereas the on-policy
arm is higher by $+0.090$ and $+0.098$ (128/128 both). That is not evidence that
retention makes prediction easier: C ranks only blocks the policy already kept, so
an evicted block can never be missed, and the value neither isolates a feedback
effect nor measures the coverage lost to eviction. Neither contrast reproduces the
archived 0.948-to-0.643 drop of \S\ref{sec:historical}, which also differs in
corpus, aggregation and numerics; attribution would need a label-only arm on
identical trajectories, candidates, features and scorer, so that gap stays
unexplained.

\subsection{Selector semantics across the three evidence objects}
\label{sec:protocol}
Table~\ref{tab:protocol} lists the rules that decide what each object measures. Two
differ on purpose: the masked loop (pinned-rerun rules) recomputes its kept set every
eight steps, so an
excluded block can return, and budgets $\mathrm{round}(0.2n)$ blocks with $n$
counting the blocks of generated tokens too, whereas the reference deletes
irreversibly and fixes its budget from the prompt.

\begin{table}[t]
\centering
\caption{Selector and label semantics. Blocks are 32 tokens; $n$ is the candidate
count, $L$ the prompt length, $K$ the oracle size of Exp\#7. The loop's budget wins when
it is below its mandatory set (7 blocks on one prompt per model).}
\label{tab:protocol}
\footnotesize
\setlength{\tabcolsep}{2.5pt}
\begin{tabular}{@{}>{\raggedright\arraybackslash}p{0.17\columnwidth}>{\raggedright\arraybackslash}p{0.26\columnwidth}>{\raggedright\arraybackslash}p{0.26\columnwidth}>{\raggedright\arraybackslash}p{0.22\columnwidth}@{}}
\toprule
rule & v2 collector & masked loop & physical reference \\
\midrule
attention & head/token mean; ragged block over members & head-group/query mean; ragged block padded & head/query/token mean \\
EMA & decay 0.9 from first decode value & H2O-style: last-32 sum; reconstructed: last-8 EMA, 0.9 & decay 0.9; last-64 prefill queries \\
within view & $\bar a/(\max\bar a+10^{-9})$ & $\bar a/(\max\bar a+10^{-9})$ & $\bar a/\max\bar a$ \\
cross view & previous-layer EMA top-$\lceil0.1n\rceil$; layer 0: own EMA & same rank rule; layer 0: own EMA & same rank rule; layer 0: zero \\
query / age & post-RoPE $q$; age from block creation & query 0.5; age since last top-$k$ & unused \\
label & future top-$\lceil0.1n\rceil$, real lookahead & current unmasked top-$K$ (Exp\#7) & unused \\
budget & --- & round$(0.2n)$, every 8 steps & $\lfloor0.2\lceil L/32\rceil\rfloor$, prompt-fixed \\
mandatory & --- & reconstructed: 4 sink~\cite{streamingllm} + 4 recent; H2O-style: half heavy, half recent & 1 sink + 1 recent \\
re-admission, new tokens & --- & allowed at each decision; new tokens unmasked, budget grows & irreversible; new tokens eligible, fixed budget \\
\bottomrule
\end{tabular}
\end{table}

\subsection{Thresholds and online updates}
\label{sec:online}\label{sec:conformal}
A score can drive a fixed-size top-$k$ set or a thresholded set
$S_{\ell,t}=\{b:f(\mathbf x_{\ell,b,t})\ge\tau_\ell\}$. The former fixes logical
retention, whereas the latter lets set size vary. We evaluate static and online
variants separately. Threshold results are reported in \S\ref{sec:guardkv};
its empirical in-distribution performance is not a hard per-request risk guarantee.
